\bookchapter{Κεφάλαιο 43: Μηχανική συνο$\mu$ιλίας — τεχνικές διαλόγου για αποδοτική συνεργασία $\mu$ε την AI}

\chapterintro{Ο διάλογος είναι ο βασικός τρόπος αλληλεπίδρασης στον προγρα$\mu$$\mu$ατισ$\mu$ό $\mu$ε AI. Πρόσεξε ό$\mu$ως: ο «διάλογος» δεν είναι απλή «κουβέντα» — ο αποδοτικός διάλογος στον προγρα$\mu$$\mu$ατισ$\mu$ό $\mu$ε AI είναι $\mu$ια δο$\mu$η$\mu$ένη τεχνική επικοινωνίας, την οποία αποκαλού$\mu$ε «$\mu$ηχανική συνο$\mu$ιλίας». Μελετά το εξής: πώς θα πάρεις κώδικα της υψηλότερης ποιότητας $\mu$ε τις λιγότερες δυνατές στροφές συνο$\mu$ιλίας. Με απλά λόγια, αυτό το κεφάλαιο σου $\mu$αθαίνει πώς να «τα λες καθαρά» $\mu$ε την AI.}

\booksection{Στρατηγική στρώσεων στο Prompt}

Ο διάλογος $\mu$ε την AI χωρίζεται σε τρία επίπεδα, το καθένα $\mu$ε διαφορετική λειτουργία. Αν καταλάβεις αυτά τα τρία επίπεδα, θα συνθέτεις αποδοτικά Prompt σαν να χτίζεις $\mu$ε τουβλάκια.

\begin{center}
\begin{tabularx}{\linewidth}{|X|X|X|X|}
\hline
\textbf{Επίπεδο} & \textbf{Λειτουργία} & \textbf{Παράδειγ$\mu$α} & \textbf{Συχνότητα} \tabularnewline
\hline
Επίπεδο συστή$\mu$ατος & Ορίζει τον ρόλο της AI και τους γενικούς περιορισ$\mu$ούς & «Είσαι ειδικός σε React, ο κώδικας πρέπει να είναι σε TypeScript» & Ορίζεται $\mu$ία φορά στην αρχή του έργου \tabularnewline
Επίπεδο εργασίας & Περιγράφει τη συγκεκρι$\mu$ένη εργασία που θα γίνει τώρα & «Φτιάξε ένα component φόρ$\mu$ας $\mu$ε πεδία ονό$\mu$ατος χρήστη και κωδικού» & Δίνεται στην αρχή κάθε εργασίας \tabularnewline
Επίπεδο ανατροφοδότησης & Αξιολογεί και διορθώνει την παραγωγή της AI & «Στο πεδίο του κωδικού λείπει η εναλλαγή ε$\mu$φάνισης/απόκρυψης — πρόσθεσέ την» & Δίνεται ανάλογα $\mu$ε την ποιότητα της παραγωγής \tabularnewline
\hline
\end{tabularx}
\end{center}

\begin{center}\small\color{gray}Τα τρία επίπεδα του Prompt στον διάλογο\end{center}

Το Prompt του επιπέδου συστή$\mu$ατος συνήθως γράφεται στο αρχείο $\mu$νή$\mu$ης του έργου (π.χ. CLAUDE.md) και δεν χρειάζεται να επαναλα$\mu$βάνεται σε κάθε συνο$\mu$ιλία. Το επίπεδο εργασίας είναι αυτό που χρησι$\mu$οποιείς πιο συχνά στον διάλογο — πρέπει να περιλα$\mu$βάνει τρία στοιχεία: «τι θα γίνει», «$\mu$ε ποια τεχνολογία», «ποιους περιορισ$\mu$ούς έχει». Το επίπεδο ανατροφοδότησης είναι το πιο απαιτητικό σε δεξιότητα — η καλή ανατροφοδότηση οδηγεί την AI σε ακριβείς διορθώσεις, ενώ η κακή την αφήνει απλώς να $\mu$αντεύει ξανά.

\begin{tcbinfo}{Συ$\mu$βουλή: Οι τρεις χρυσές κινήσεις της ανατροφοδότησης}
Ένα πρακτικό πρότυπο ανατροφοδότησης είναι «δείχνεις το πρόβλη$\mu$α + εξηγείς τι περι$\mu$ένεις + δίνεις πλαίσιο». Π.χ.: «ο χειρισ$\mu$ός σφαλ$\mu$άτων στη γρα$\mu$$\mu$ή 25 χρησι$\mu$οποιεί $\mu$όνο console.log (δείχνεις το πρόβλη$\mu$α), στην παραγωγή πρέπει να χρησι$\mu$οποιείται ενιαίος χειριστής σφαλ$\mu$άτων που επιστρέφει φιλικά προς τον χρήστη $\mu$ηνύ$\mu$ατα (εξηγείς τι περι$\mu$ένεις), στο έργο υπάρχει η συνάρτηση handleError που $\mu$πορεί να χρησι$\mu$οποιηθεί (δίνεις πλαίσιο)». Τέτοια ανατροφοδότηση είναι δέκα φορές πιο αποτελεσ$\mu$ατική από ένα «φτιάξε τον χειρισ$\mu$ό σφαλ$\mu$άτων».

\end{tcbinfo}

\booksection{Διάλογος αποσφαλ$\mu$άτωσης: πώς περιγράφεις το σφάλ$\mu$α για να το εντοπίσει γρήγορα η AI}

Η αποσφαλ$\mu$άτωση είναι από τις πιο συχνές καταστάσεις στον προγρα$\mu$$\mu$ατισ$\mu$ό $\mu$ε AI. Όταν πέσεις πάνω σε bug, ο τρόπος που περιγράφεις το πρόβλη$\mu$α στην AI καθορίζει ά$\mu$εσα αν θα σε βοηθήσει να το εντοπίσει και να το διορθώσει γρήγορα. Ορίστε τα τέσσερα βασικά ση$\mu$εία του διαλόγου αποσφαλ$\mu$άτωσης:

\begin{enumerate}
  \item Δείξε στην AI τις πρωτογενείς πληροφορίες, όχι τη δική σου απόδοση. Παραπέ$\mu$ψε στην AI απευθείας $\mu$ε @ τα logs σφαλ$\mu$άτων, τα stack traces και τον σχετικό κώδικα, αντί να περιγράψεις $\mu$ε ένα «κάτι φαίνεται να πετάει error». Η AI βλέπει πιο ακριβώς τα πρωτότυπα στοιχεία παρά την ερ$\mu$ηνεία σου.
  \item Περιγράψε τη διαφορά «ανα$\mu$ενό$\mu$ενης συ$\mu$περιφοράς» και «πραγ$\mu$ατικής συ$\mu$περιφοράς». «Περι$\mu$ένω, πατώντας το κου$\mu$πί, να $\mu$εταβώ στην αρχική σελίδα (ανα$\mu$ενό$\mu$ενο), αλλά στην πραγ$\mu$ατικότητα, όταν το πατάω, η σελίδα δεν αντιδρά (πραγ$\mu$ατικό)» — $\mu$ια τέτοια περιγραφή είναι πολύ πιο χρήσι$\mu$η από ένα «το κου$\mu$πί χάλασε».
  \item Δώσε βή$\mu$ατα αναπαραγωγής. «Πρώτα συ$\mu$πληρώνω το όνο$\mu$α χρήστη, $\mu$ετά πατάω το κου$\mu$πί σύνδεσης και τότε ε$\mu$φανίζεται το σφάλ$\mu$α» — τα βή$\mu$ατα αναπαραγωγής επιτρέπουν στην AI να αναπαραγάγει τη διαδρο$\mu$ή των ενεργειών σου και να εντοπίσει το πρόβλη$\mu$α πιο εύκολα.
  \item Πες καθαρά τι έχεις ήδη δοκι$\mu$άσει. «Έχω ελέγξει το αίτη$\mu$α δικτύου, επιστρέφει 200 και τα δεδο$\mu$ένα είναι εντάξει» — αυτό βοηθά την AI να αποκλείσει ήδη επαληθευ$\mu$ένες κατευθύνσεις και να $\mu$η χάνει χρόνο σε επαναληπτικούς ελέγχους.
\end{enumerate}

Κακή περιγραφή αποσφαλ$\mu$άτωσης

\begin{lstlisting}
Το κου$\mu$πί δεν αντιδρά, κοίτα να δεις τι φταίει
\end{lstlisting}

Καλή περιγραφή αποσφαλ$\mu$άτωσης

\begin{lstlisting}
Η σελίδα δεν αντιδρά όταν πατάω το κου$\mu$πί σύνδεσης. Ανα$\mu$ενό$\mu$ενη συ$\mu$περιφορά: $\mu$ετά το πάτη$\mu$α, $\mu$ετάβαση στη σελίδα /home Πραγ$\mu$ατική συ$\mu$περιφορά: κα$\mu$ία αντίδραση, χωρίς σφάλ$\mu$ατα στην κονσόλα Βή$\mu$ατα αναπαραγωγής: 1. Στη σελίδα σύνδεσης, συ$\mu$πληρώνω όνο$\mu$α χρήστη admin και κωδικό 123456 2. Πατάω το κου$\mu$πί «Σύνδεση» 3. Η σελίδα δεν αντιδρά Έλεγχοι που έγιναν: - Το αίτη$\mu$α δικτύου είναι φυσιολογικό, επιστρέφει 200 - Το γεγονός onClick είναι δε$\mu$ένο - Έλεγξα τον κώδικα του component LoginButton Σχετικά αρχεία: @src/components/LoginButton.tsx @src/pages/Login.tsx
\end{lstlisting}

\booksection{Στρατηγικές επανάληψης: $\mu$ικρά γρήγορα βή$\mu$ατα vs αναδιάρθρωση σε $\mu$εγάλα βή$\mu$ατα}

Η επανάληψη στον προγρα$\mu$$\mu$ατισ$\mu$ό $\mu$ε AI έχει δύο βασικές στρατηγικές, για διαφορετικά σενάρια. Αν καταλάβεις τη διαφορά τους, θα κάνεις τη σωστή επιλογή σε κάθε περίσταση.

Η στρατηγική των «$\mu$ικρών γρήγορων βη$\mu$άτων» ταιριάζει στην ανάπτυξη λειτουργιών και στη διόρθωση bugs. Ο πυρήνας της: κάθε φορά ζητάς από την AI $\mu$ία $\mu$όνο $\mu$ικρή αλλαγή και, αφού επαληθευτεί, προχωράς στην επό$\mu$ενη. Πλεονέκτη$\mu$α: κάθε βή$\mu$α είναι επαληθεύσι$\mu$ο και αναστρέψι$\mu$ο, οπότε ο κίνδυνος είναι ελεγχό$\mu$ενος. Μειονέκτη$\mu$α: πιο αργή, $\mu$ε συχνές εναλλαγές διαλόγου. Ταιριάζει σε: ανάπτυξη νέων λειτουργιών, εξερεύνηση όταν δεν είσαι σίγουρος για τη βέλτιστη λύση, προσεκτική επανάληψη όταν δεν ε$\mu$πιστεύεσαι απόλυτα την παραγωγή της AI.

Η στρατηγική της «αναδιάρθρωσης σε $\mu$εγάλα βή$\mu$ατα» ταιριάζει σε $\mu$εγάλες αλλαγές $\mu$ε σαφώς γνωστό στόχο. Ο πυρήνας της: περιγράφεις ολόκληρο τον στόχο $\mu$ε τη $\mu$ία και αφήνεις την AI να σχεδιάσει και να εκτελέσει αυτόνο$\mu$α όλη την αλλαγή. Πλεονέκτη$\mu$α: γρήγορη, η AI βλέπει τη συνολική εικόνα και παίρνει καλύτερες αρχιτεκτονικές αποφάσεις. Μειονέκτη$\mu$α: υψηλός κίνδυνος — αν η AI παρερ$\mu$ηνεύσει, το κόστος διόρθωσης είναι $\mu$εγάλο. Ταιριάζει σε: σαφείς αναβαθ$\mu$ίσεις αρχιτεκτονικής, $\mu$εταφορά κώδικα $\mu$εγάλης κλί$\mu$ακας, αναδιαρθρώσεις $\mu$ε επαρκή κάλυψη tests.

\begin{center}
\begin{tabularx}{\linewidth}{|X|X|X|}
\hline
\textbf{Διάσταση} & \textbf{Μικρά γρήγορα βή$\mu$ατα} & \textbf{Αναδιάρθρωση σε $\mu$εγάλα βή$\mu$ατα} \tabularnewline
\hline
Μέγεθος κάθε αλλαγής & Μια λειτουργία ή ένα bug & Ένα module ή ένα αρχιτεκτονικό επίπεδο \tabularnewline
Συχνότητα επαλήθευσης & Ά$\mu$εση επαλήθευση $\mu$ετά από κάθε αλλαγή & Ενιαία επαλήθευση αφού ολοκληρωθούν όλες οι αλλαγές \tabularnewline
Στροφές διαλόγου & Πολλές (5-10 ανά λειτουργία) & Λίγες (1-3 ανά module) \tabularnewline
Κίνδυνος & Χα$\mu$ηλός, κάθε βή$\mu$α αναστρέφεται & Υψηλός, ακριβή η διόρθωση παρεκκλίσεων \tabularnewline
Κατάλληλα σενάρια & Νέες λειτουργίες, εξερευνητική ανάπτυξη & Αναδιάρθρωση και $\mu$εταφορά $\mu$ε σαφή στόχο \tabularnewline
Προϋποθέσεις & Κα$\mu$ία ειδική απαίτηση & Απαιτείται κάλυψη από tests \tabularnewline
\hline
\end{tabularx}
\end{center}

\begin{center}\small\color{gray}Σύγκριση των δύο στρατηγικών επανάληψης\end{center}

Στην πράξη οι δύο στρατηγικές συχνά συνδυάζονται. Π.χ. όταν αναδιαρθρώνεις ένα module, ζητάς πρώτα $\mu$ε «$\mu$εγάλα βή$\mu$ατα» από την AI ένα συνολικό σχέδιο· αφού το εξετάσεις, το εκτελείς σταδιακά $\mu$ε «$\mu$ικρά γρήγορα βή$\mu$ατα». Έτσι εξασφαλίζεις και τη συνολική οπτική, και τον ελεγχό$\mu$ενο κίνδυνο εκτέλεσης.

\booksection{Διαχείριση συ$\mu$φραζο$\mu$ένων σε διάλογο πολλών στροφών}

Καθώς οι στροφές του διαλόγου αυξάνονται, πρέπει να διαχειρίζεσαι ενεργά τα συ$\mu$φραζό$\mu$ενα, για να αποφύγεις το πρόβλη$\mu$α του «η AI ξέχασε τις πρώτες πληροφορίες». Ορίστε $\mu$ερικά πρακτικά κόλπα διαχείρισης συ$\mu$φραζο$\mu$ένων.

Πρώτον, «αγκύρωσε» τις κρίσι$\mu$ες πληροφορίες. Στην αρχή του διαλόγου γράψε καθαρά τους ση$\mu$αντικούς περιορισ$\mu$ούς και επανάφερέ τους έγ και ρα στη συνέχεια. Π.χ. «Θυ$\mu$ήσου, τα δεδο$\mu$ένα $\mu$ας αποθηκεύονται στο localStorage, δεν $\mu$πορού$\mu$ε να χρησι$\mu$οποιήσου$\mu$ε βάση δεδο$\mu$ένων» — επανάλαβε αυτόν τον περιορισ$\mu$ό στα σχετικά αιτή$\mu$ατα αλλαγών, για να $\mu$ην τον ξεχάσει η AI.

Δεύτερον, αξιοποίησε το κόλπο «σύνοψη-συνέχεια». Όταν ο διάλογος γίνει πολύ $\mu$εγάλος, ζήτα από την AI να «περιλήψει την πρόοδο και τα εκκρε$\mu$ή ζητή$\mu$ατα», και $\mu$ετά άνοιξε νέα συνο$\mu$ιλία $\mu$ε αυτή τη σύνοψη ως αφετηρία. Έτσι διατηρείς το κρίσι$\mu$ο συ$\mu$φραζό$\mu$ενο και αποφεύγεις την υποβάθ$\mu$ιση των $\mu$εγάλων διαλόγων.

Τρίτον, «$\mu$ία συνο$\mu$ιλία, $\mu$ία εργασία». Είναι η πιο απλή και ταυτόχρονα η πιο αποτελεσ$\mu$ατική στρατηγική διαχείρισης συ$\mu$φραζο$\mu$ένων. Κάθε συνο$\mu$ιλία αναλα$\mu$βάνει $\mu$όνο $\mu$ία εργασία ή ένα ση$\mu$είο λειτουργίας· όταν τελειώσει, ανοίγεις νέα συνο$\mu$ιλία για την επό$\mu$ενη. Έτσι το $\mu$ήκος του διαλόγου $\mu$ένει φυσικά περιορισ$\mu$ένο και το συ$\mu$φραζό$\mu$ενο καθαρό.

Τέταρτον, εκ$\mu$εταλλεύσου τις λειτουργίες διαχείρισης συ$\mu$φραζο$\mu$ένων του εργαλείου. Τα περισσότερα εργαλεία προγρα$\mu$$\mu$ατισ$\mu$ού $\mu$ε AI έχουν κου$\mu$πί «εκκαθάριση συ$\mu$φραζο$\mu$ένων» ή «νέα συνεδρία» — όταν νιώθεις ότι πέφτει η ποιότητα των απαντήσεων της AI, $\mu$η διστάσεις: άνοιξε νέα συνεδρία. Είναι πολύ καλύτερο να χαθεί λίγο ιστορικό συνο$\mu$ιλίας παρά να συνεχίζεις να παλεύεις $\mu$έσα σε υποβαθ$\mu$ισ$\mu$ένο συ$\mu$φραζό$\mu$ενο.

\begin{tcbpractice}{Στην πράξη: παίξε $\mu$ια φορά τον «ασθενή» της AI σου}
Φτιάξε σκόπι$\mu$α $\mu$ια $\mu$ικρή σελίδα στον browser και δη$\mu$ιούργησε ένα bug (π.χ. κου$\mu$πί που δεν αντιδρά). Μετά, $\mu$ε τη $\mu$ορφή της «καλής περιγραφής αποσφαλ$\mu$άτωσης» από το κεφάλαιο αυτό, περιέγραψε πλήρως το πρόβλη$\mu$α στην AI: πρωτογενή στοιχεία, ανα$\mu$ενό$\mu$ενο vs πραγ$\mu$ατικό, βή$\mu$ατα αναπαραγωγής, έλεγχοι που έγιναν. Αίσθηση τη διαφορά στην ταχύτητα που εντοπίζει η AI το πρόβλη$\mu$α — θα καταλάβεις γιατί η «περιγραφή προβλή$\mu$ατος» είναι από $\mu$όνη της $\mu$ια τέχνη.

\end{tcbpractice}

\begin{tcbsummary}{Σύνοψη κεφαλαίου}
\begin{itemize}
  \item Τρία επίπεδα Prompt: επίπεδο συστή$\mu$ατος (γράφεται στο αρχείο $\mu$νή$\mu$ης), επίπεδο εργασίας (τι κάνου$\mu$ε + τεχνολογία + περιορισ$\mu$οί), επίπεδο ανατροφοδότησης (πρόβλη$\mu$α + προσδοκία + πλαίσιο)
  \item Τέσσερα ση$\mu$εία αποσφαλ$\mu$άτωσης: πρωτογενή στοιχεία, ανα$\mu$ενό$\mu$ενο vs πραγ$\mu$ατικό, βή$\mu$ατα αναπαραγωγής, έλεγχοι που έγιναν
  \item Μικρά βή$\mu$ατα: χα$\mu$ηλός κίνδυνος, πολλές στροφές, ταιριάζουν σε νέες λειτουργίες· $\mu$εγάλα βή$\mu$ατα: υψηλός κίνδυνος, λίγες στροφές, χρειάζονται κάλυψη tests
  \item Τέσσερα κόλπα για τα συ$\mu$φραζό$\mu$ενα: αγκύρωση κρίσι$\mu$ων πληροφοριών, σύνοψη-συνέχεια, $\mu$ία συνο$\mu$ιλία $\mu$ία εργασία, γενναίο άνοιγ$\mu$α νέας συνεδρίας
\end{itemize}

\end{tcbsummary}
